\chapter{Structure Sets the Spectrum}
\label{chapter:act-I}

% Chapter 4 opening:

Chapter~\ref{chapter:prelims} divided the spectrum of a connectivity matrix into
two parts with different roles. The leading eigenvalue sets the gain, and with it
the dynamical regime the network occupies; it is the part we rescale and
therefore the part we control. The eigenvalues beneath it are where a particular
wiring leaves its mark. That division was introduced as a way of reading a
spectrum. This chapter measures it, on a human connectome and on a family of
substrates that destroy parts of its organisation one at a time. Nothing here is
simulated. Every quantity reported is computed from the connectivity matrix
alone, before any input is presented, which is what makes the spectrum the first
link in the chain and the only one that can be examined without running the
network at all.

What separates a connectome from a randomised network of the same size is not
the bulk. Measured in the raw units of the weights, the bulk edges of the
connectome and of every substrate tested against it agree to within a few per
cent. What separates them is the outlier. The connectome's leading eigenvalue
stands more than three times clear of its own bulk edge, where every other
substrate stands under two, and essentially the whole difference between them
sits in that one number. This is worth stating carefully, because the convention
in this literature is to divide each matrix by its own leading eigenvalue before
comparing, and in those units the connectome's bulk is the narrow one. That
narrowness is a property of the divisor rather than of the bulk, and the divisor
is the one quantity the substrates differ in. Stated in the direction the
measurement supports, the connectome's weight assignment buys a large spectral
gap rather than a compact bulk.

The assignment is the operative object, and not the graph or the weights taken
separately. A control that holds the connectome's own edge set and its own
weights and reshuffles only which edge carries which value lands with the
randomised substrates rather than with the connectome. No measure of the binary
graph distinguishes the two, and no statistic of the weight values distinguishes
them either, since the multiset is identical. Neither is sufficient on its own.
What produces the gap is the correspondence between them.

That localises the effect without explaining it. Which property of the
assignment is responsible remains open, and the one candidate specific enough to
test here, a weighted rich club in which the heaviest edges run between the
highest-degree parcels, is ruled out by the measurement. A spectrum also says
nothing about what a network does with what it is given. It describes an
operator, and where each part of that operator ends up in the activity the
operator generates is the question of Chapter~\ref{chapter:act-II}. Two things 
leave this chapter for it. The first is a scale that every substrate carries with 
it, the point at which its own bulk edge reaches one; it makes a second way of 
matching substrates possible, and the closing section sets out why neither way is 
neutral. The second is a mode standing clear of the bulk, handed on as an object 
to be decomposed rather than as an explanation.






\section{The substrate and its null ladder}
\label{sec:act1-substrate}
\bigskip

The substrate is a human structural connectome at 448 cortical
parcels, % METH A.4
built for this thesis as a group consensus over the diffusion-weighted scans of
a publicly released healthy-adult cohort \cite{griffa2019structural,
suarez2021learning} rather than taken from a published matrix. Two of its
properties carry through the whole thesis. It is symmetric, and therefore
normal, which fixes the form its spectrum can take. And it is non-negative,
since tractography counts streamlines and cannot represent an inhibitory
connection. The same build at
$N = 1000$ % METH A.4
supports the change of scale in Section~\ref{sec:act1-scale}. Construction is in
Chapter~\ref{chapter:methods}.

Against the connectome stand three randomised substrates, each destroying some
part of its organisation and holding the rest fixed.
Table~\ref{tab:act1-preservation} sets out what each keeps, and
Figure~\ref{fig:F19} shows what that looks like as a graph. An Erdős–Rényi
graph matches the node count, edge count and density exactly rather than in
expectation, and panel~(d) is the result: connections scattered uniformly, with
none of the block structure of panel~(a). A degree-preserving rewire adds the
degree sequence, so it carries the connectome's heavy-tailed degree
distribution intact, visible in the marginal above panel~(c), while the blocks
go the same way. Weights are painted onto both by drawing, with replacement,
from the connectome's own multiset of edge
weights, % METH B.0
so the ladder varies where weight sits without changing what weights are
available.

The third randomised substrate is not a null but a control. The weight-permuted
variant takes the connectome's own graph and its own weights and permutes the
weights across the edges, and panels~(a) and~(b) of Figure~\ref{fig:F19} are the
same image because of it: same edge set, same degrees, and the same weight
multiset exactly rather than in distribution, since the values are reordered
rather than resampled. % METH B.1
The magnified block below makes the one difference legible. In (i) and (j) the
same cells are filled and the colours in them are not the same colours. No
measure of the binary graph separates the two substrates, which is what
Table~\ref{tab:act1-topology} reports: identical clustering, identical
modularity on a fixed partition, identical assortativity and identical
efficiency, as any such measure must return on one graph. Against the
connectome the control isolates weight assignment, holding topology and weight
statistics fixed together; against the degree-preserving rewire it isolates
topology, since the two draw their weights from the same pool and differ in the
graph. Each of the three randomised substrates is generated at ten
seeds; % METH B.0
the connectome is one fixed graph and returns the same matrix every time. 
Three further substrates are not among the four drawn here and are described 
in Chapter~\ref{chapter:methods}; Figure~\ref{fig:S4} draws the whole family
together and their spectra are in Appendix~\ref{app:wider-ladder}.


% Null model preservation table, four drawn rungs only.
\begin{table}[ht]
    \centering
    \caption[What each substrate preserves]{\textbf{What each substrate
    preserves and what it destroys.} The four substrates the figures of this
    chapter draw, ordered by how much of the connectome each keeps, most at the
    top. All four share the connectome's node count, edge count, density and
    non-negative weights, so those properties are not columns here; what
    separates them is everything below that. The weight-permuted variant is a
    control rather than a null, preserving the graph entire and destroying only
    the assignment of weights to edges. Generation procedures, and the three
    further substrates this chapter does not draw, are in
    Chapter~\ref{chapter:methods}.}
    \label{tab:act1-preservation}
    \begin{tabular}{lcccc}
        \toprule
        & degree & clustering & modularity & weight \\
        & sequence &  &  & assignment \\
        \midrule
        \textbf{connectome} & \checkmark & \checkmark & \checkmark & \checkmark \\
        weight-permuted & \checkmark & \checkmark & \checkmark & --- \\
        degree rewire & \checkmark & --- & --- & --- \\
        Erdős–Rényi & --- & --- & --- & --- \\
        \bottomrule
    \end{tabular}
\end{table}

% Four Substrates Adjacency Matrices Figure
\begin{figure}[htbp]
    \centering
    \includegraphics[width=\textwidth]{F19.png}
    \caption[The four substrates, as graphs]{%
    \textbf{No measure of the binary graph separates the connectome from its
    control.} The four substrates of this chapter, drawn under one node
    ordering: hemisphere, then community, with communities detected once on the
    connectome and reused for every panel. All four hold
    448 % METH A.4
    nodes and 5{,}323 % METH A.6
    undirected edges at a density of 5.32\% % METH A.6
    by construction. \textbf{(a--d)} Binary adjacency, with the degree marginal
    above each panel. Panels (a) and (b) are the same image: the
    weight-permuted control preserves the connectome's edge set exactly, so the
    two are one graph. The degree-preserving rewire (c) keeps the marginal and
    loses the block structure; Erdős–Rényi (d) keeps neither. \textbf{(e--h)}
    Weighted adjacency on a shared logarithmic colour scale, same ordering.
    \textbf{(i--l)} One diagonal block magnified, marked on (e--h) and drawn on
    the same colour scale: one hemisphere's share of one community, the densest
    such block of between 15 and 40 nodes. % TIER0 3.13
    In (i) and (j) the same cells are
    filled, because the edge set is the same edge set; what the permutation
    changes is which weight sits in each of them, drawn from the whole graph's
    multiset rather than from this block's. The same nodes in (k) and (l) form
    no block at all. Randomised substrates are drawn at the seed used
    throughout this chapter. Topology statistics are in
    Table~\ref{tab:act1-topology}; where these substrates differ is the
    spectrum, in Section~\ref{sec:act1-spectrum}.}
    \label{fig:F19}
\end{figure}

% Topology comparison table
\begin{table}[ht]
    \centering
    \caption[Topology of the four substrates]{\textbf{Topology of the four
    drawn substrates.} All four statistics are measures of the binary graph,
    weights discarded; medians over ten seeds for the two randomised
    substrates, which are single values for the connectome and its control.
    Modularity is evaluated on one partition, detected once on the connectome
    and applied unchanged to all four, so the four scores are comparable;
    detecting communities separately would find them in Erdős–Rényi too. The
    connectome and the weight-permuted control return identical values on every
    statistic, as any measure of the binary graph must, since they are the same
    graph. Rows are ordered as in Table~\ref{tab:act1-preservation}.}
    \label{tab:act1-topology}
    \begin{tabular}{lrrrr}
        \toprule
        & mean & modularity & degree & global \\
        substrate & clustering & (fixed part.) & assortativity & efficiency \\
        \midrule
        \textbf{connectome} & 0.4277 & 0.5486 & 0.1067 & 0.4064 \\ % TIER0 3.13
        weight-permuted & 0.4277 & 0.5486 & 0.1067 & 0.4064 \\ % TIER0 3.13
        degree-preserving rewire & 0.0697 & $-0.0002$ & $-0.0192$ & 0.4789 \\ % TIER0 3.13
        Erdős–Rényi & 0.0529 & $-0.0036$ & $-0.0068$ & 0.4820 \\ % TIER0 3.13
        \bottomrule
    \end{tabular}
\end{table}



\section{An outlier and a bulk}
\label{sec:act1-spectrum}
\bigskip
% Outline items 3 and 4. Carries A1.1, A1.2, A1.3. Figures F1, F2.

The connectome's spectrum is a dense band about the origin and one eigenvalue
standing well clear of it at the right-hand edge, which is panel~(a) of
Figure~\ref{fig:F1}. It is also entirely real, as symmetry requires, and not
merely to within a tolerance one could work with: across all forty ladder cells
at $N = 448$ the largest imaginary part in the stored spectra is exactly
zero, % A1 2.5
and the symmetry check holds in every row. % A1 2.5
Figure~\ref{fig:F1} is therefore drawn on the real line rather than in the
complex plane, where a second dimension would carry nothing.

The separation itself is guaranteed rather than an accident of this particular
graph. Non-negativity gives $W$ a real, positive leading eigenvalue, largest in
modulus and with a non-negative eigenvector, and the only question a given
substrate answers is how far clear of everything else that eigenvalue sits.
Three terms follow. The \emph{Perron root} $\lambda_1$ is that leading
eigenvalue, the \emph{bulk} is the band holding the remainder, and the
\emph{gap} is the separation between them. Measuring the gap needs an edge for
the bulk, and the one used throughout is $\mathrm{bulk95}$, the 95th percentile
of $|\lambda|$ over the spectrum as a fraction of
$\lambda_1$. % A1 2.6 finding 1
All three are properties of a single substrate, measured before any comparison
is made.

\begin{figure}[htbp]
    \centering
    \includegraphics[width=\textwidth]{F1.png}
    \caption[The connectome's weight assignment buys a large spectral gap]{%
    \textbf{The connectome's weight assignment buys a large spectral gap.} Eigenvalues of the un-rescaled recurrent matrix $W$ on the
    all-positive human substrate ($N = 448$ % TIER0 2.1
    self-built consensus, empirical weights).
    \textbf{(a--d)} One substrate per row, raw units, shared axis: the bulk
    shaded to $\pm$ the absolute bulk radius, the gap out to $\pm\lambda_1$ in
    grey. Sighting down the column, the four bulk edges hold while $\lambda_1$
    (dotted) separates and the grey gap grows from row d to row a; the
    weight-permuted control (b) sits with the nulls. Each row draws one
    spectrum, for the nulls the seed whose $\mathrm{bulk95}$ is nearest the
    median, so all four carry equal sampling noise. \textbf{(e)} The same
    spectra as cumulative distributions, each divided by its own $\lambda_1$,
    the normalisation the connectome-reservoir literature uses, and every nominal 
    spectral-radius-matched comparison in this thesis with it.
    The curves separate everywhere; the 95\% crossings are
    $\mathrm{bulk95}$ itself. \textbf{(f)} The gap ratio per seed, with
    $\mathrm{bulk95}$ on the right-hand axis. The connectome's ten seeds
    coincide exactly.}
    \label{fig:F1}
\end{figure}

The connectome's Perron root stands
more than three % (3.078) TIER0 3.1
times clear of its own bulk edge at $N = 448$. % TIER0 3.1
Every other substrate the figures draw stands between
1.81 % TIER0 3.1
and 1.92 % TIER0 3.1
times clear of its own. That ratio, the Perron root over the absolute bulk
radius, is the gap ratio, and the connectome's is at least
1.6-fold % TIER0 3.1(b)
above every rung of the four-substrate ladder it is measured against. The gap is
not bought by a narrower bulk. It is bought by a higher outlier over a bulk that
is essentially everyone's.

Figure~\ref{fig:F1} draws the eigenvalues of the un-rescaled matrix $W$, one
substrate per row, on a shared axis in the raw units of the weights, and the
first thing to do with it is to sight down the column. The four bulk bands sit
at the same place. Their edges fall between
0.0587 % TIER0 3.1
and 0.0614, % TIER0 3.1
a spread of 4.4\% % TIER0 3.1
of their mean at $N = 448$. % TIER0 3.1
What moves is the outlier. The Perron root runs from
0.1061 % TIER0 3.1
at the Erdős–Rényi rung to 0.1889 % TIER0 3.1
at the connectome, so the grey region separating the bulk from $\lambda_1$ grows
steadily up the column while the coloured band beneath it stays put.
Figure~\ref{fig:F2} makes the same decomposition as medians across the ladder,
the absolute bulk in panel~(a) and the Perron root in panel~(b), and
Table~\ref{tab:act1-ladder} gives the values at both parcellations.

Reported instead as a fraction of each substrate's own Perron root, those same
four bulk radii spread 47.3\%. % TIER0 3.1
Nothing about the bulks has moved between those two numbers. Only the divisor
has. Panel~(e) is where the divisor becomes visible: it redraws the same spectra
as cumulative distributions after each substrate has been divided by its own
$|\lambda_1|$, the normalisation every spectral-radius-matched comparison in the
connectome-reservoir literature
uses \cite{suarez2021learning, suarez2024connectome, milisav2026neuromorphic}.
In those units the curves separate everywhere and the connectome's reads as the
narrow one, with $\mathrm{bulk95}$ legible at the crossings of the 95\% level.
Read that way, the connectome has an anomalously compact bulk.
Read in the units of the weights it has nothing of the kind: the four bulks
agree, and what the connectome has instead is a Perron root
1.78 % TIER0 3.1
times the Erdős–Rényi rung's. % TIER0 3.1(a)
The two halves of the figure are the same spectra one division apart, and the
compactness belongs to the divisor. Stated in the direction the measurement
supports, the difference is an anomalously large spectral gap.

The ratio has two other names, and all three are the same number:
$|\lambda_1| / \mathrm{abs\_bulk} = 1/\mathrm{bulk95} = \mathrm{sr\_crit}$. % TIER0 3.1
The gap ratio, the inverse bulk and the critical scale are one quantity, which
is why panel~(f) can carry $\mathrm{bulk95}$ on its right-hand axis, and why the
name $\mathrm{sr\_crit}$ is the one worth adopting: it names the quantity after
what it measures, the nominal spectral radius at which the median seed's bulk
edge reaches one.\footnote{The edge in question is the 95th percentile of
$|\lambda|$, % A1 2.6 finding 1
so at $\sigma = \mathrm{sr\_crit}$ five per cent % A1 2.6 finding 1
of modes already exceed unity. $\mathrm{sr\_crit}$ is a criticality scale each
substrate carries with it, not the point at which that substrate becomes
supercritical. Chapter~\ref{chapter:methods} gives the aggregation the identity
holds under.}

\begin{table}[ht]
    \centering
    \caption[The substrate ladder at both scales]{\textbf{The substrate ladder
    at both scales.} Medians over ten seeds; the connectome is one fixed graph,
    so its seeds coincide exactly and only the nulls are resampled.
    $|\lambda_1|$ and the absolute bulk radius are in raw units of $W$; the
    absolute bulk is the product of the medians,
    $\mathrm{median}(\mathrm{bulk95}) \times \mathrm{median}(|\lambda_1|)$, and
    $\mathrm{bulk95}$ is the ratio of the two. Spread is the range over the
    mean. The entire between-variant difference sits in $|\lambda_1|$, which is
    why the two spreads differ by close to an order of magnitude. 
    Rows are ordered as in Table~\ref{tab:act1-preservation} rather than by value; 
    the gap-ratio column is not monotone at $N = 1000$.}
    \label{tab:act1-ladder}
    \begin{tabular}{lrrrr}
        \toprule
        substrate & absolute bulk & $|\lambda_1|$ & $\mathrm{bulk95}$ & gap ratio \\
        \midrule
        \multicolumn{5}{l}{\emph{$N = 448$}} \\
        connectome & 0.0614 & 0.1889 & 0.3249 & \textbf{3.078} \\ % TIER0 3.1
        weight-permuted & 0.0599 & 0.1152 & 0.5203 & 1.922 \\ % TIER0 3.1
        degree-preserving rewire & 0.0595 & 0.1115 & 0.5338 & 1.873 \\ % TIER0 3.1
        Erdős–Rényi & 0.0587 & 0.1061 & 0.5535 & 1.807 \\ % TIER0 3.1
        \addlinespace
        spread (range / mean) & 4.4\% & --- & 47.3\% & --- \\ % TIER0 3.1
        \midrule
        \multicolumn{5}{l}{\emph{$N = 1000$}} \\
        connectome & 0.0585 & 0.2333 & 0.2509 & \textbf{3.985} \\ % TIER0 3.1(a)
        weight-permuted & 0.0584 & 0.1398 & 0.4176 & 2.395 \\ % TIER0 3.1(a)
        degree-preserving rewire & 0.0590 & 0.1358 & 0.4346 & 2.301 \\ % TIER0 3.1(a)
        Erdős–Rényi & 0.0553 & 0.1348 & 0.4102 & 2.438 \\ % TIER0 3.1(a)
        \addlinespace
        spread (range / mean) & 6.4\% & --- & 48.5\% & --- \\ % TIER0 3.1
        \bottomrule
    \end{tabular}
\end{table}


%%%%%%%%%%%%%%%%%%%%%%%%%%%%%%%%%%%%%%%%%%%%


\section{Weight assignment and the gap}
\label{sec:act1-assignment}
\bigskip
% Outline item 5. Carries A1.5. Figure F1 (row b), Table 4.3.

What produces the gap is settled by the second row of Figure~\ref{fig:F1}. The
weight-permuted control differs from the connectome in nothing but which edge
carries which weight, and it lands with the nulls, at a gap ratio of
1.92 % TIER0 3.1
against the connectome's 3.08. % TIER0 3.1
The graph is held entire and the weight multiset is held exactly, so neither is
sufficient on its own. What buys the gap is the correspondence between them.

Two further rungs corroborate this without carrying it. Each holds more of the
connectome's topology fixed than the degree-preserving rewire does, one
preserving its community structure and one its clustering, and neither recovers
the gap; their gap ratios at both parcellations, and the margin over all six, 
are in Appendix~\ref{app:wider-ladder}. Each of those rungs fixes one
topological statistic. The control fixes the entire graph, which is why the
contrast that settles the question is the one against it.

The contrast localises the effect to the assignment. It does not say which
property of the assignment is responsible, and no mechanism is proposed for why
the assignment raises the Perron root. One account is specific enough to test on
these substrates. In a weighted rich club the heavy edges run between the
high-degree nodes, so they share endpoints and reinforce one another in the
leading eigenvalue, and a weight permutation destroys exactly that arrangement.
The prediction was registered before the measurement and it fails: edge weight
and endpoint degree product correlate at
$-0.124$ % TIER0 3.14
in the connectome, negative rather than positive and with a bootstrap interval
excluding zero, while the control sits at zero as its construction
requires. % TIER0 3.14
The connectome's heavy edges run between its lower-degree parcels, not between
its hubs, so whatever the assignment is doing to the Perron root, it is not
assembling a heavily weighted core. Appendix~\ref{app:mechanism} gives the test,
its two other predictions, and the account left standing by it, which is not
tested here.


%%%%%%%%%%%%%%%%%%%%%%%%%%%%%%%%%%%%%%%%%%%%


\section{The gap across scale}
\label{sec:act1-scale}
\bigskip
% Outline item 6. Carries A1.4. Figures F2 (both scales), S1 (appendix).

The connectome's spectral gap survives a more than two-fold
% (2.2-fold) TIER0 1.4
change in the number of nodes. The whole ladder was rebuilt at
$N = 1000$, % TIER0 2.1
and panel~(c) of Figure~\ref{fig:F2} is where that shows: the connectome's gap
ratio is 3.99 % TIER0 3.1
against 2.30 to 2.44 % TIER0 3.1
for the three rungs drawn beside it, again at least
1.63-fold % TIER0 3.1(b)
clear of every one, with the full ladder at both parcellations in
Table~\ref{tab:act1-ladder}.
Two parcellations are one step and there is no third, so
this is a robustness check across a single change of scale rather than a
demonstration of invariance.

What does not survive is the near-identity of the bulks. At $N = 448$ the four
absolute bulk radii agree to 4.4\% % TIER0 3.1
of their mean; at $N = 1000$ they spread 6.4\%, % TIER0 3.1
or 6.9\% % TIER0 3.1
if the per-seed products are taken before the median rather than after.
\emph{Near-identical} describes $N = 448$; at $N = 1000$ the right word is
\emph{close}. The loosening is narrower than the summary figure suggests, since
three of the four substrates still sit within
1.1\% % TIER0 3.1(a)
of one another and the whole of the difference is the Erdős–Rényi rung sitting
low. Scale also unsettles the ordering among the nulls, which does not hold
between the two parcellations on either the bulk edge or the absolute
bulk.\footnote{The largest absolute bulk belongs to the connectome at
$N = 448$ and to the degree-preserving rewire at
$N = 1000$. % TIER0 3.1(a)
Figure~\ref{fig:S1} rebuilds Figure~\ref{fig:F1} at the larger parcellation with
its rows in ladder order, so the gap ratios run down that column out of
sequence.} What is scale-robust is the connectome's separation from the other
three, not the coincidence of the bulks and not any arrangement among the nulls.


\begin{figure}[htbp]
    \centering
    \includegraphics[width=\textwidth]{F2.png}
    \caption[The difference is a gap, and the gap is what carries across
    scale]{%
    \textbf{The difference between the connectome and its nulls is a spectral
    gap, and what carries across scale is the gap.} The four drawn substrates at
    both parcellations, ten seeds each (solid $N = 448$, % TIER0 2.1
    translucent $N = 1000$). % TIER0 2.1
    All three panels are zero-based and use one aggregation throughout, the
    median over seeds of $\mathrm{bulk95}$ and of $|\lambda_1|$, so (c) is
    exactly (b) divided by (a). \textbf{(a)} The absolute bulk radius,
    near-identical across the four at $N = 448$ and merely close at $N = 1000$.
    \textbf{(b)} $|\lambda_1|$ in raw units, where the substrates do separate.
    \textbf{(c)} The gap ratio, identically
    $1/\mathrm{median}(\mathrm{bulk95})$ and identically the critical scale
    $\mathrm{sr\_crit}$. Across a 2.2-fold % TIER0 1.4
    change in $N$ the connectome's separation from the other three holds while
    the near-identity in (a) loosens, so what is scale-robust is the separation
    and not the coincidence of the bulks.}
    \label{fig:F2}
\end{figure}




\section{What carries forward}
\label{sec:act1-handover}
\bigskip
% Outline item 7. Carries A1.6. No new measurements. Figure F3 is in
% Chapter 10; this section states the argument, not the demonstration.

The first thing to leave this chapter is $\mathrm{sr\_crit}$, the nominal
spectral radius at which a substrate's own bulk edge reaches one. Because the
gap ratio is also the inverse bulk, every substrate arrives carrying one, and
the four drawn here do not carry the same one. That is a consequence of the
result rather than a further finding, and it is what makes a second way of
matching substrates possible at all.

Comparing substrates requires putting them at the same operating point, and on
these substrates there are two ways to do it. Scaling each matrix by its own
$|\lambda_1|$ to a nominal spectral radius $\sigma$ sets the leading eigenvalue
to $\sigma$ exactly and leaves the bulk wherever it falls; this is the
normalisation the connectome-reservoir literature
uses \cite{suarez2021learning, suarez2024connectome, milisav2026neuromorphic},
and it is the one panel~(e) of Figure~\ref{fig:F1} draws. The alternative is the
product $\sigma \cdot \mathrm{bulk95}$, which holds the bulk edge fixed instead,
so that at $\sigma \cdot \mathrm{bulk95} = 1$ every substrate's bulk edge
reaches one whatever its Perron root. The second axis exists only because each
substrate carries its own $\mathrm{sr\_crit}$, which is what this chapter
supplies.

Neither axis is neutral, and the second is not a repair of the first. Each holds
fixed exactly what the other lets vary, and since these substrates differ in
essentially one quantity, no third axis holds both at once. A comparison on the
nominal axis is a comparison at matched Perron roots and unmatched bulks; one on
the matched-bulk axis is the reverse. Two consequences follow, and both are
enforced throughout what follows. Every result is quoted with the axis it was
measured on, and results on the two axes are never set against each other as
though they disagreed, because they are measurements of different things. And
the two axes cover different ranges, since holding the bulk fixed compresses the
reachable operating points by each substrate's own $\mathrm{bulk95}$: a sweep
that runs to $\sigma = 8$ on the nominal axis reaches much less far in the
matched coordinate, so a comparison there is bounded by where the substrates'
ranges still overlap, and a feature at the edge of that overlap must not be read
as though the sweep continued past it. Section~\ref{sec:methods-matching} in Methods 
shows what the choice costs on a real result.

One use of $\mathrm{sr\_crit}$ is worth flagging now, because it is a choice
rather than a consequence. The filter that separates supercritical from
subcritical operating points in Chapters~\ref{chapter:act-II}
and~\ref{chapter:act-III} is the connectome's own critical scale, applied to
every substrate rather than each substrate's own. It is the conservative
option, since it samples each null further above its own critical point than the
null's own scale would, and it is set out again where it is used.

The second thing handed on is the Perron mode. It stands clear of the bulk, and
the connectome's weight assignment is what puts it there. What the mode does is
another question, and not one a spectrum can answer: a spectrum describes an
operator, while what a network computes depends on where each part of that
operator ends up in the activity the operator generates. That is what
Chapter~\ref{chapter:act-II} measures, and the mode arrives there as an object
to be decomposed rather than as an explanation.



